\section{Protection of vulnerable people: a duty of care, written down}
\label{sec:protection}

The sixth appeal is protective rather than purely compassionate. Health debates
return again and again to the groups most exposed to harm: children, the elderly,
disabled patients, veterans, and rural communities. For medical AI the specific
worries are bias against minority populations, errors that fall hardest on the
vulnerable, and consent that is assumed rather than confirmed.

[DRAFTING INSTRUCTIONS: in the full-narrative, anchor the protection appeal in
\texttt{review/references/narrative\_refs.bib} entries
\texttt{obermeyer2019dissecting} and \texttt{fda2021samd} (the FDA AI/ML SaMD
action plan) and in the author's site-documentation and CFR adaption works
(\texttt{Kawchak2026PhysicalAIOncologyClinical},
\texttt{Kawchak2026Adaption21CFRPartb}, on 21 CFR Part 50 informed consent).
Reproduce Mermaid figure
\texttt{review/mermaid/diagrams/15-vulnerable-population-safeguards.md} as colored
TikZ. Add a body-width table listing each protected group, its specific risk, and
the bill safeguard that addresses it.]

\begin{narrativefig}
\node[nfone] (a) at (0,0) {Subgroup data};
\node[nfdec] (b) at (3.8,0) {Consent +\\bias audit};
\node[nfact] (c) at (7.8,0) {Reported safeguard};
\draw[nfedge] (a)--(b); \draw[nfedge] (b)--(c);
\end{narrativefig}
\vspace{-0.2cm}
\figcaption{Figure 7. Protection written into the procedure (placeholder;
expanded from Mermaid 15 in the full-narrative).}

Protection that is written into statute, measured, and reported is stronger than
protection that depends on the good intentions of whoever happens to be on duty.
